\documentclass[11pt]{article}
\usepackage[utf8]{inputenc}
\usepackage{amsmath,amssymb,amsthm}
\usepackage{enumerate}
\usepackage[margin=1in]{geometry}
\usepackage{hyperref}

\newtheorem{theorem}{Theorem}
\newtheorem{lemma}[theorem]{Lemma}
\newtheorem{corollary}[theorem]{Corollary}
\newtheorem{definition}[theorem]{Definition}
\newtheorem{remark}[theorem]{Remark}
\newtheorem{proposition}[theorem]{Proposition}

\newcommand{\lf}{\mathsf{lf}}
\newcommand{\nd}{\mathsf{nd}}
\newcommand{\cas}{\mathsf{cas}}
\newcommand{\rec}{\mathsf{rec}}
\newcommand{\niter}{\mathsf{niter}}
\newcommand{\eval}{\mathsf{eval}}
\newcommand{\reify}{\mathsf{reify}}
\newcommand{\codeTerm}{\mathsf{codeTerm}}
\newcommand{\codeReify}{\mathsf{codeReify}}
\newcommand{\diagCode}{\mathsf{diagCode}}
\newcommand{\diag}{\mathsf{diag}}
\newcommand{\thS}{\mathsf{thS}}
\newcommand{\thTerm}{\mathsf{thTerm}}
\newcommand{\eqTree}{\mathsf{eqTree}}
\newcommand{\transCode}{\mathsf{transCode}}
\newcommand{\swapCode}{\mathsf{swapCode}}
\newcommand{\Tree}{\mathsf{Tree}}
\newcommand{\Term}{\mathsf{Term}}
\newcommand{\Eq}{\mathsf{Eq}}
\newcommand{\dThS}{D_{\thS}}
\newcommand{\ValidProof}{\mathsf{ValidProof}}
\newcommand{\ThSResult}{\mathsf{ThSResult}}
\newcommand{\ProvEq}{\mathsf{ProvEq}}
\newcommand{\coreTree}{\mathsf{coreTree}}
\newcommand{\CoreInv}{\mathsf{CoreInv}}
\newcommand{\comp}{\mathsf{comp}}
\newcommand{\matchSub}{\mathsf{matchSub}}

\title{G\"odel's Second Incompleteness Theorem \`a la Rose:\\
  Plan for the Binary Tree Calculus}
\author{}
\date{March 2026}

\begin{document}
\maketitle

\section{Rose's Proof Structure}

Rose~\cite{rose1961} proves consistency and G\"odel~II for Ternary
Recursive Arithmetic~$T$ entirely within the object system, without
introducing a formula layer, provability predicate, or Hilbert--Bernays
derivability conditions.  The argument has five components:

\begin{enumerate}
\item \textbf{Theorem enumerator} $th(y)$: the $y$-th theorem of~$T$,
  defined by cases on $y \bmod 6$ (one case per axiom schema plus
  transitivity).

\item \textbf{Soundness} (Theorem~10): $vl(x, th(y)) = 0$ for all
  valuations~$x$ and all~$y$.  Every enumerated equation is true.

\item \textbf{Companion function} $e(y)$: given the number of an
  equation $A = B$, returns the number of $A \mathbin{+} B$
  (the successor-paired companion).  The companion is always false:
  $vl(x, e(y)) = 1$ whenever $vl(x, y) = 0$ (Theorem~11).

\item \textbf{Consistency} (Theorem~13): $th(y) \neq e(th(z))$ for
  all $y, z$.  A theorem is never a companion of a theorem.  Immediate
  from soundness + companion falsity.

\item \textbf{Representability and G\"odel~II} (Theorems~14--18):
  the substitution function $se$ and numeral function $nu$ are
  \emph{representable} in~$T$ (Theorem~17).  Combined with the
  diagonal construction, the consistency sentence
  $th(y) \neq e(th(z))$ is shown to be undecidable in~$T$.
\end{enumerate}

\noindent
No formula type, no implication, no D2/D3, no L\"ob argument.
The entire proof lives in the equational calculus.


\section{Our Setting: Binary Tree Calculus}

We work over binary trees ($\lf$, $\nd\;a\;b$) with intrinsically
scoped terms ($\Term\;n$), primitive recursion ($\rec$), and
right-spine iteration ($\niter$).  The existing formalization provides:

\begin{center}
\begin{tabular}{lll}
\textbf{Rose} & \textbf{Ours} & \textbf{Status} \\
\hline
$th(y)$ & $\thS : \Tree \to \Tree$ & Done (\texttt{ThInt.agda}) \\
$vl(x, th(y)) = 0$ & $\thS$-valid $: \ValidProof\;p \to \ThSResult\;(\thS\;p)$
  & Done \\
$e(y)$ & companion function $\comp$ & \textbf{Not yet} \\
$th(y) \neq e(th(z))$ & consistency via $\comp$ & \textbf{Not yet} \\
Thm~14--16 (D1 internal) & $\dThS$ + equational D1 & Partial \\
Thm~17 (representability) & $\mathsf{chiFormTerm}$-correct & Partial \\
Thm~18 (G\"odel~II) & self-referential undecidability & \textbf{Not yet} \\
\end{tabular}
\end{center}

\noindent
The existing \texttt{GodelII.agda} follows Guard/Hilbert--Bernays
(formula layer + L\"ob argument), not Rose.  The goal is now to
follow Rose's object-level approach.


\section{The Companion Function}\label{sec:companion}

Rose's $e(y)$ takes an equation code and produces a provably false
companion.  In our setting, an equation code is $\nd\;L\;R$
where $L$ and $R$ are term codes.  We need a computable function
$\comp : \Tree \to \Tree$ such that:

\begin{enumerate}
\item[(C1)] $\comp$ is total and computable (definable as a $\Term\;1$).
\item[(C2)] If $\ThSResult\;(\nd\;L\;R)$ (i.e., $L = \codeTerm\;l$,
  $R = \codeTerm\;r$, $\eval\;l = \eval\;r$), then $\comp\;(\nd\;L\;R)$
  is a well-formed equation code $\nd\;L'\;R'$ with $\eval\;l' \neq \eval\;r'$.
\end{enumerate}

\begin{definition}[Companion function]
The simplest candidate:
\[
  \comp\;(\nd\;L\;R) = \nd\;L\;(\nd\;L\;R)
\]
interpreting: the LHS stays the same, the RHS becomes $\nd\;L\;R$
(the pair of the original two sides).  Since $\eval\;l$ is a single
tree and $\nd\;(\eval\;l)\;(\eval\;r)$ is always $\nd$-headed,
$\eval\;l \neq \nd\;(\eval\;l)\;(\eval\;r)$ whenever $\eval\;l = \lf$
or more generally when $\eval\;l$ has a different top-level structure.
\end{definition}

\begin{remark}
The simplest \emph{universal} companion: $\comp\;(\nd\;L\;R) =
\nd\;(\codeTerm\;\mathsf{leaf})\;(\codeTerm\;(\mathsf{pair}\;\mathsf{leaf}\;\mathsf{leaf}))$.
This is the equation $\lf = \nd\;\lf\;\lf$, which is always false.
But this is constant---it doesn't depend on the input.  Rose's $e(y)$
depends on the input because $A \mathbin{+} B$ pairs the two sides,
making the companion \emph{structurally related} to the theorem.
For our purposes, even a constant companion suffices for consistency
(Theorem~13 analogue), but a dependent companion may be needed for
the self-referential argument (Theorem~18).
\end{remark}

A more faithful analogue of Rose's $e$: define $\comp$ so that
the companion of $\nd\;L\;R$ is an equation whose LHS \emph{is}
the original equation code (as a tree), and whose RHS is something
different.  For instance:
\[
  \comp\;(\nd\;L\;R) = \nd\;(\codeReify\;(\nd\;L\;R))\;(\codeReify\;\lf)
\]
This encodes ``the equation code itself equals $\lf$'', which is
always false (since $\nd\;L\;R \neq \lf$).  Moreover, both sides
are valid term codes ($\codeReify$ always produces $\codeTerm\;(\reify\;\cdot)$),
so the companion is a well-formed equation code.


\section{Consistency \`a la Rose}\label{sec:consistency}

\begin{theorem}[Rose's Theorem~13 analogue]
For all $y, z : \Tree$:
\[
  \thS\;y \neq \comp\;(\thS\;z)
\]
\end{theorem}

\begin{proof}[Proof strategy]
By $\ThSResult$, $\thS\;z = \nd\;(\codeTerm\;l)\;(\codeTerm\;r)$
with $\eval\;l = \eval\;r$ (for valid~$z$; for all~$z$ via $\coreTree$).
The companion $\comp\;(\thS\;z)$ is an equation code whose two sides
evaluate to \emph{different} trees (by construction of $\comp$).
By $\ThSResult$ applied to~$y$, $\thS\;y$ always encodes a
\emph{true} equation.  A true equation cannot equal a false equation.
\end{proof}

\noindent
This is exactly Rose's argument: soundness says $\thS$ only produces
true equations; the companion produces false equations; therefore a
theorem is never a companion of a theorem.

\begin{remark}
We already have the pieces:
\begin{itemize}
  \item $\thS$-valid gives $\ThSResult\;(\thS\;y)$ for valid proofs.
  \item $\mathsf{falseEq\text{-}unprovable}$ in \texttt{Godel.agda}
    already proves: if $\eval\;s \neq \eval\;t$, then
    $\thS\;y \neq \nd\;(\codeTerm\;s)\;(\codeTerm\;t)$ for valid~$y$.
  \item $\mathsf{falseEq\text{-}never}$ strengthens this to all~$y$
    (not just valid proofs) via $\coreTree$.
\end{itemize}
The companion function just packages these into Rose's form.
\end{remark}


\section{Representability (Rose's Theorems~14--17)}\label{sec:represent}

This is the hard part.  Rose proves that the substitution function
$se(x, y, z)$ and numeral function $nu(x)$ are \emph{representable}:
there exist terms in~$T$ that compute these functions, and the
equations witnessing this are provable in~$T$ (i.e., appear as
$th(y)$ outputs).

In our setting, representability means: for each primitive recursive
function $f$ used in the diagonal construction, there exists a
$\Term\;n$ computing~$f$, and the equation
$\codeTerm\;(\text{term for } f) = \codeTerm\;(\reify\;(f\;\vec{a}))$
is in the range of~$\thS$.

\subsection{What we already have}

\begin{itemize}
\item $\thTerm : \Term\;1$ with $\thTerm$-correct: $\eval_y\;\thTerm = \thS\;y$.
  This \emph{internalizes} $\thS$.
\item $\dThS$: from proof tree~$p$, produces a proof of the reflexivity
  of $\thS\;p$.  This is Rose's Theorem~14 restricted to reflexivity.
\item $\mathsf{chiFormTerm}$: encodes formulas as closed terms with
  $\eval\;(\mathsf{chiFormTerm}\;A) = \mathsf{chiForm}\;A$.
\item $\matchSub$-correct: $\matchSub$ computes $\eqTree$ internally.
\item $\codeTerm$/$\codeReify$ with round-trip properties.
\end{itemize}

\subsection{What we need}

\begin{enumerate}
\item \textbf{Internal $\codeTerm$}: a $\Term\;1$ that computes
  $\codeTerm\;(\reify\;x)$ from~$x$.  This is $\codeReify$ internalized.
  We already have $\codeReify : \Tree \to \Tree$ at the metalevel
  and $\codeReify$-correct linking it to $\codeTerm\;(\reify\;\cdot)$.
  The internal version is a $\rec$-based term; the unit tests in
  \texttt{ThInt.agda} (\texttt{crBase}, \texttt{crStep}) suggest
  this is already partially built.

\item \textbf{Internal $\comp$}: a $\Term\;1$ computing the companion.
  If $\comp$ is simple (e.g., $\comp\;(\nd\;L\;R) =
  \nd\;(\codeReify\;(\nd\;L\;R))\;(\codeReify\;\lf)$), this is
  just $\mathsf{pair}\;(\text{codeReify-term applied to } v_0)\;
  (\text{codeReify-term applied to } \lf)$.

\item \textbf{Representability of $\comp \circ \thS$}: the equation
  $\thS(y) \neq \comp(\thS(z))$ must be expressible as a $\thS$ output
  for the self-referential argument.  This requires showing that
  the \emph{equation} ``$\comp(\thS(z))$ is not a theorem'' can be
  diagonalized.

\item \textbf{The diagonal argument}: using the existing
  $\diag$/$\diagCode$ machinery, construct a sentence~$G$ such that
  $\eval\;G = \eqTree\;(\thS\;c_G)\;(\comp\;(\thS\;c_G))$ where
  $c_G$ is the code of~$G$'s schema.  If $G$ is provable ($\thS\;y =
  \text{code of } G$), then by soundness $\eval\;G = \lf$, meaning
  $\thS\;c_G = \comp\;(\thS\;c_G)$, contradicting consistency.
  If $\neg G$ is provable, similar contradiction.  Therefore $G$ is
  undecidable.
\end{enumerate}


\section{Concrete Plan for the Next Session}\label{sec:plan}

\subsection{Step 1: The companion function (small)}

Define $\comp : \Tree \to \Tree$ and $\mathsf{compTerm} : \Term\;1$.
Prove:
\begin{itemize}
  \item $\comp$ applied to any $\ThSResult$ output gives a false equation.
  \item $\eval_y\;\mathsf{compTerm} = \comp\;y$.
\end{itemize}
This directly gives Rose's Theorem~13 analogue:
$\thS\;y \neq \comp\;(\thS\;z)$ for all valid~$y, z$.

\subsection{Step 2: Consistency via companion (small)}

Combine Step~1 with $\mathsf{falseEq\text{-}unprovable}$ or
$\mathsf{falseEq\text{-}never}$ from \texttt{Godel.agda}. This
gives an object-level consistency theorem without any formula layer.

\subsection{Step 3: Internal codeReify (medium)}

Build a $\Term\;1$ that computes $\codeReify\;x$ from~$x$.
The infrastructure (\texttt{crBase}, \texttt{crStep} in
\texttt{ThInt.agda}) suggests this is a $\rec$-based term.
Prove correctness.

\subsection{Step 4: The G\"odel sentence via companion (medium)}

Using the existing diagonal machinery ($\diag$, $\diagCode$,
$\mathsf{godelSchema}$), construct a G\"odel sentence~$G$ whose
evaluation checks whether $\thS$ at the self-referential input
equals the companion of $\thS$ at the same input.  Formally:
\[
  \eval\;G = \eqTree\;(\thS\;c)\;(\comp\;(\thS\;c))
\]
where $c = \codeTerm\;(\mathsf{godelSchema}\;\ldots)$.

By $\eqTree$-$\thS$-$\lf$ and the companion construction, $\eval\;G$
will be $\mathsf{falseT}$ (since $\thS\;c$ is always $\nd$-headed
and $\comp$ produces something different).  This gives a true
sentence.  By the companion-consistency theorem, the equation
$G = \reify\;\mathsf{falseT}$ is unprovable.

\subsection{Step 5: Representability and undecidability (hard)}

Prove that the companion-of-$\thS$ is representable: there exists
a proof tree whose $\thS$ output encodes the equation
$\mathsf{compTerm}\;(\thTerm\;y) = \comp\;(\thS\;y)$.
Combined with the diagonal, this gives Rose's Theorem~18:
the consistency sentence is undecidable.

This is the hardest step and may require extending $\thS$ with
additional computation axioms (cas-pair, rec-pair) so that
$\thS$ can prove representability equations.


\section{What the Existing GodelII.agda Provides}

The existing \texttt{GodelII.agda} (600 lines) is not wasted work.
It provides:

\begin{itemize}
\item The L\"ob-style G\"odel~II as a \emph{sanity check}: the
  abstract argument works, confirming the proof system is set up
  correctly.
\item $\ProvEq$ with closure under sym/trans/axioms: directly
  useful for showing equations are in the range of~$\thS$.
\item The chi-translation ($\mathsf{chiForm}$, $\mathsf{chiFormTerm}$):
  infrastructure for encoding truth values as terms, which the
  companion function and representability proofs will use.
\item $\mathsf{chiProvAt}$-to-$\ProvEq$: connects chi-checks to
  provability witnesses.
\end{itemize}

\noindent
The new Rose-style development will be a \emph{separate section}
of \texttt{GodelII.agda} (or a new file \texttt{Rose/GodelIIRose.agda}),
building on the same infrastructure but following Rose's proof
structure.


\begin{thebibliography}{9}
\bibitem{rose1961}
H.\,E.\ Rose,
\emph{On the consistency and undecidability of recursive arithmetic},
Zeitschrift f\"ur mathematische Logik und Grundlagen der Mathematik
\textbf{7} (1961), 124--135.

\bibitem{guard1963}
J.\,R.\ Guard,
\emph{Lecture notes on recursive arithmetic},
Technical report, Princeton University, 1963.
\end{thebibliography}

\end{document}
